%! Sine and Cosine Function Values — Unit 3, Lesson 3.3 — Student Packet Cover Sheet
\documentclass[10pt]{article}
\usepackage{precalculus-article}
\usepackage{precalculus-boxes}
\usepackage{ltablex}
\keepXColumns

\begin{document}

% Full-bleed navy banner
\begin{tikzpicture}[remember picture, overlay]
  \fill[navy] (current page.north west)
    rectangle ([yshift=-0.9in]current page.north east);
\end{tikzpicture}
\vspace{-0.8in}

\begin{center}
  {\color{white}\textbf{\LARGE Precalculus}}\\[4pt]
  {\color{skymid}\large\textbf{Unit 3: Trigonometric and Polar Functions}}\\[6pt]
  {\color{white}\textbf{\large Lesson 3.3 \quad Sine and Cosine Function Values}}
\end{center}

\vspace{0.22in}
\namedateperiod
\vspace{0.18in}

\begin{learningtargetbox}
\small
\begin{enumerate}[leftmargin=1.5em, itemsep=5pt]
  \item I can find the coordinates of a point on a circle of radius $r$ for an angle $\theta$ using $P = (r\cos\theta,\, r\sin\theta)$.
  \item I can use 45-45-90 and 30-60-90 triangle geometry to determine exact sine and cosine values at multiples of $\pi/6$ and $\pi/4$.
  \item I can apply quadrant sign conventions to obtain correct exact values when the angle lies outside the first quadrant.
\end{enumerate}
\end{learningtargetbox}

\vspace{0.14in}

\begin{tocbox}
\small
\renewcommand{\arraystretch}{1.5}
\begin{tabularx}{\linewidth}{@{} c l X r @{}}
\rowcolor{sky}
\textbf{\#} & \textbf{Component} & \textbf{Description} & \textbf{Score} \\
\midrule
1 & Warm-Up        & Pythagorean theorem, special right triangles, unit-circle basics & \blank{1.2cm} \\
2 & Guided Notes   & Special angle exact values; $P=(r\cos\theta, r\sin\theta)$        & \blank{1.2cm} \\
3 & Group Activity & Tiered exact-value and coordinate problems                         & \blank{1.2cm} \\
4 & Exit Ticket    & Quadrant-II/III exact values; radius-$r$ coordinate problem        & \blank{1.2cm} \\
5 & Homework       & Fluency with QII--QIV exact values; AP-style MC; symmetry proof   & \blank{1.2cm} \\
\midrule
  & \multicolumn{2}{r}{\textbf{Total}} & \blank{1.2cm} \\
\end{tabularx}
\end{tocbox}

\end{document}
